\section{Fiche : machines frigorifiques / pompes à chaleur (ex. R134a)}

\subsection{Principe et schéma}

Un frigo et une pompe à chaleur (PAC) sont la \textbf{même machine} : un \textbf{cycle récepteur} à deux
sources (il faut fournir du travail $W>0$ au fluide). Seule change l'énergie « utile » :
\begin{itemize}[leftmargin=1.4em]
\item \textbf{Frigo} : utile $=Q_2$ (chaleur \emph{prise} à la source froide, l'intérieur du frigo).
\item \textbf{PAC} : utile $=-Q_1$ (chaleur \emph{cédée} à la source chaude, l'intérieur de la maison) ;
la chaleur prise à la source froide (air extérieur) est considérée \textbf{gratuite}.
\end{itemize}

\begin{center}
\begin{tikzpicture}[scale=0.8, >=Stealth, thick]
% boucle
\draw (0,0) -- (0,3) -- (4,3) -- (4,0) -- cycle;
% evaporateur
\draw (0,1.1) -- (0.35,1.4) -- (-0.35,1.7) -- (0,2.0);
\node at (-1.35,1.95) {\small\shortstack{Évapo-\\rateur}};
% condenseur
\draw (4,1.1) -- (4.35,1.4) -- (3.65,1.7) -- (4,2.0);
\node at (5.35,1.95) {\small\shortstack{Conden-\\seur}};
% compresseur
\draw (1.5,3) -- (2,2.35) -- (2.5,3);
\node at (2,3.5) {\small Compresseur};
% detendeur
\draw (1.5,0.4) -- (2.5,0.4) -- (2,0) -- cycle;
\draw (1.5,-0.4) -- (2.5,-0.4) -- (2,0) -- cycle;
\node at (2,-0.75) {\small Détendeur};
% sources
\node at (-3.1,1.9) {\small\shortstack{Source\\froide}};
\draw[->] (-2.5,1.5) -- (-2.0,1.15);
\node at (7.1,1.9) {\small\shortstack{Source\\chaude}};
\draw[->] (6.5,1.5) -- (6.0,1.15);
% points
\node at (0.3,2.15) {\small 1}; \node at (2.35,2.65) {\small 2};
\node at (3.7,2.15) {\small 3}; \node at (1.65,-0.15) {\small 4};
\end{tikzpicture}
\end{center}

\begin{center}
\renewcommand{\arraystretch}{1.5}
\begin{tabularx}{\linewidth}{@{}l l X@{}}
\toprule
\textbf{Étape} & \textbf{Transfo} & \textbf{Ce qui se passe} \\
\midrule
1 $\to$ 2 & Compression isentropique & vapeur BP $\to$ vapeur HP, $T$ et $p$ augmentent \\
2 $\to$ 3 & Condensation isobare (HP) & le fluide cède $Q_1<0$ à la source chaude, se liquéfie \\
3 $\to$ 4 & Détente isenthalpique ($h$=cste) & dans la soupape, $p$ et $T$ chutent, $W_{ut}=0$ et $Q=0$ \\
4 $\to$ 1 & Évaporation isobare (BP) & le fluide reçoit $Q_2>0$ de la source froide, se vaporise \\
\bottomrule
\end{tabularx}
\end{center}

\subsection{Le diagramme des frigoristes $(\log p,\,h)$}

En pratique on place le cycle sur un diagramme $\log p$ - $h$ du fluide (fourni à l'examen). Points \emph{réels}
notés $1',2',3',4'$ : on prévoit une \textbf{surchauffe} ($\sim5\,^\circ$C au-dessus de $T_{\text{évap}}$ en sortie
évaporateur, pour garantir un gaz sec avant le compresseur) et un \textbf{sous-refroidissement}
($\sim5\,^\circ$C sous $T_{cond}$ en sortie condenseur, pour garantir du liquide pur avant la soupape).

\begin{center}
\begin{tikzpicture}[scale=1, >=Stealth]
\draw[->] (0,0) -- (8.5,0) node[right] {$h$ [kJ/kg]};
\draw[->] (0,0) -- (0,5.3) node[above] {$\log p$};
% dome
\draw[thick] (1.2,0.4) to[out=70,in=180] (3.1,4.4) to[out=0,in=110] (5.4,0.4);
% pressure levels
\draw[dashed] (0,3.6) -- (7.6,3.6) node[right]{\small $p_{cond}$};
\draw[dashed] (0,1.1) -- (6.8,1.1) node[right]{\small $p_{\text{évap}}$};
% points
\coordinate (p4) at (2.55,1.1); \coordinate (p1) at (4.0,1.1);
\coordinate (p2) at (5.85,3.6); \coordinate (p3) at (2.85,3.6);
\draw[thick,->] (p1) to[out=95,in=-100] (p2); % compression isentropique
\draw[thick,->] (p2) -- (p3); % condensation
\draw[thick,->] (p3) -- (p4); % détente isenthalpique
\draw[thick,->] (p4) -- (p1); % évaporation
\foreach \p/\lab/\pos in {p1/1'/below,p2/2'/above,p3/3'/above,p4/4'/below}
  {\fill (\p) circle (1.6pt) node[\pos] {\lab};}
\end{tikzpicture}
\end{center}

\begin{methode}
\textbf{Méthode pas à pas pour placer le cycle et lire les enthalpies}
\begin{enumerate}[leftmargin=1.6em]
\item Repérer sur l'axe vertical les deux pressions $p_{\text{évap}}$ (à $T_{\text{évap}}$) et $p_{cond}$ (à $T_{cond}$)
grâce à la cloche de saturation (intersection avec la courbe = pression de saturation à cette température).
\item Placer \textbf{1'} sur l'isobare $p_{\text{évap}}$, à $T_{\text{évap}}+$surchauffe (à droite de la cloche, en zone
vapeur).
\item Placer \textbf{2'} sur l'isobare $p_{cond}$, en suivant \textbf{l'isentropique} passant par 1'
(ou via $\eta_{isos}$ si compression non idéale, voir plus bas).
\item Placer \textbf{3'} sur l'isobare $p_{cond}$, à $T_{cond}-$sous-refroidissement (à gauche de la cloche,
zone liquide).
\item Placer \textbf{4'} \textbf{à la verticale de 3'} (même $h$, détente isenthalpique) sur l'isobare
$p_{\text{évap}}$.
\item Lire $h_{1'},h_{2'},h_{3'},h_{4'}$ directement sur l'axe des $h$.
\end{enumerate}
\end{methode}

\subsection{Formules}

\begin{center}
\renewcommand{\arraystretch}{1.7}
\begin{tabularx}{\linewidth}{@{}l X@{}}
\toprule
Travail compresseur (par kg) & $w=h_{2'}-h_{1'}$ \\
Chaleur cédée (condenseur, par kg) & $q_1=h_{3'}-h_{2'}$ \ (négatif) \\
Chaleur reçue (évaporateur, par kg) & $q_2=h_{1'}-h_{4'}$ \ (positif) \\
Vérification & $q_1+q_2+w=0$ \\
$PSN$ (frigo) & $PSN=\dfrac{Q_2}{W}=\dfrac{h_{1'}-h_{4'}}{h_{2'}-h_{1'}}$ \\
$COP$ (PAC) & $COP=\dfrac{-Q_1}{W}=\dfrac{h_{2'}-h_{3'}}{h_{2'}-h_{1'}}$ \\
Rendement isentropique du compresseur & $\eta_{isos}=\dfrac{h_{2isos}-h_1}{h_{2\,\text{réel}}-h_1}$ (donne
$h_{2'}$ réel à partir de $h_{2s}$ lu sur l'isentropique) \\
Bornes théoriques (Carnot, $T$ en K) & $COP_{max}=\dfrac{T_1}{T_1-T_2}$, \quad
$PSN_{max}=\dfrac{T_2}{T_1-T_2}$ \\
Débit de fluide frigorigène & $\dot m=\dfrac{P_{utile}}{|q_{utile}|}$ (kg/s), puis $P_{compresseur}=\dot m\cdot w$ \\
\bottomrule
\end{tabularx}
\end{center}

\begin{piege}
\textbf{$T_1$/$T_2$ dans les formules de Carnot = températures des \emph{sources}} (air extérieur, eau du
chauffage), \textbf{pas} les températures d'évaporation/condensation du fluide (qui sont décalées de
quelques degrés à cause de l'écart d'échange dans les échangeurs). Ne mélangez pas les deux dans une même
formule.
\end{piege}

\newpage
\subsection{Exemple complet résolu -- PAC au R134a pour le chauffage d'une maison}

\begin{enonce}
Une pompe à chaleur au R134a prélève de la chaleur à l'air extérieur pour chauffer l'eau d'un circuit de
radiateurs. Évaporation à $-5\,^\circ$C, condensation à $45\,^\circ$C, surchauffe et sous-refroidissement
de $5\,^\circ$C. Rendement isentropique du compresseur $\eta_{isos}=0{,}75$.

Lecture sur le diagramme $(\log p,h)$ du R134a : $h_{1'}=402{,}5$ kJ/kg ; $h_{2s}=435$ kJ/kg (point
isentropique) ; $h_{3'}=257$ kJ/kg.

La maison perd $P=12$ kW qu'il faut compenser. L'eau des radiateurs entre au condenseur à $35\,^\circ$C
et en ressort à $40\,^\circ$C ($c_{eau}=4185$ J/kg·K).

\emph{Trouver : le COP de la PAC, le débit de R134a, la puissance du compresseur, et le débit d'eau du
circuit de chauffage.}
\end{enonce}

\begin{correction}
\textbf{1) Point 2' réel (compression non isentropique)} :
\[\eta_{isos}=\frac{h_{2s}-h_{1'}}{h_{2'}-h_{1'}} \Rightarrow
h_{2'}=h_{1'}+\frac{h_{2s}-h_{1'}}{\eta_{isos}}=402{,}5+\frac{435-402{,}5}{0{,}75}=402{,}5+43{,}3
=445{,}8\text{ kJ/kg}\]

\textbf{2) Point 4'} (détente isenthalpique) : $h_{4'}=h_{3'}=257$ kJ/kg.

\textbf{3) Énergies échangées par kg de fluide} :
\[w=h_{2'}-h_{1'}=445{,}8-402{,}5=43{,}3\text{ kJ/kg}\]
\[q_1=h_{3'}-h_{2'}=257-445{,}8=-188{,}8\text{ kJ/kg} \ \ (\Rightarrow |q_1|=188{,}8\text{ kJ/kg})\]
\[q_2=h_{1'}-h_{4'}=402{,}5-257=145{,}5\text{ kJ/kg}\]
\emph{Vérification : } $q_1+q_2+w=-188{,}8+145{,}5+43{,}3\approx0$ \checkmark

\textbf{4) COP de la PAC} :
\eqbox{COP=\frac{-q_1}{w}=\frac{188{,}8}{43{,}3}=\mathbf{4{,}36}}

\emph{Comparaison Carnot (source froide $=-5\,^\circ$C$=268{,}15$K, source chaude $=45\,^\circ$C$=318{,}15$K)} :
$COP_{max}=\dfrac{318{,}15}{318{,}15-268{,}15}=6{,}36$. Notre COP réel ($4{,}36$) représente $\sim69\%$
du maximum théorique, ce qui est réaliste pour une PAC.

\textbf{5) Débit de fluide frigorigène} (le condenseur doit fournir $12$ kW) :
\eqbox{\dot m_{R134a}=\frac{P}{|q_1|}=\frac{12}{188{,}8}=\mathbf{0{,}0636\ kg/s}}

\textbf{6) Puissance du compresseur} :
\eqbox{P_{comp}=\dot m_{R134a}\cdot w=0{,}0636\times43{,}3=\mathbf{2{,}76\ kW}}

\emph{(Vérification cohérence : $P_{comp}+P_{\text{évap}}=P_{cond}$, soit $2{,}76+\dot m\cdot q_2=2{,}76+
0{,}0636\times145{,}5=2{,}76+9{,}25=12{,}0$ kW \checkmark)}

\textbf{7) Débit d'eau du circuit de chauffage} -- échangeur ouvert isobare, $\Delta H=Q$ :
\[P=\dot m_{eau}\cdot c_{eau}\cdot\Delta T \Rightarrow
\dot m_{eau}=\frac{P}{c_{eau}\cdot\Delta T}=\frac{12\,000}{4185\times5}\]
\eqbox{\dot m_{eau}\approx \mathbf{0{,}573\ kg/s}\ (\approx2{,}06\ \text{m}^3/\text{h})}
\end{correction}

\begin{piege}
La partie « échange thermique » (eau du chauffage, air d'un local, etc.) se traite \textbf{toujours} avec
$P=\dot m\, c\,\Delta T$ (ou $Q=mc\Delta T$ sans débit) -- c'est simplement le 1\textsuperscript{er} principe
en système ouvert isobare ($\Delta H=Q$) appliqué à un fluide \emph{sans} changement de phase. Le lien entre
les deux circuits (fluide frigorigène $\leftrightarrow$ eau) se fait via la \textbf{puissance échangée},
identique des deux côtés de l'échangeur.
\end{piege}
